\chapter*{Phần mở đầu}
\phantomsection
\addcontentsline{toc}{chapter}{Phần mở đầu}

\section*{1. Lý do chọn đề tài}
\phantomsection
\addcontentsline{toc}{section}{1. Lý do chọn đề tài}

Thương mại điện tử thời trang là lĩnh vực có nhu cầu cá nhân hóa cao. Khi mua quần áo trực tuyến, người dùng thường phải đánh giá sản phẩm qua ảnh mẫu, mô tả ngắn và bảng kích thước. Cách tiếp cận này chưa giải quyết tốt các câu hỏi thực tế như trang phục có hợp vóc dáng hay không, nên phối với món nào, có phù hợp với hoàn cảnh sử dụng hay không và khi mặc lên cơ thể thật sẽ có cảm giác thị giác như thế nào.

\medskip

Trong khi đó, sự phát triển của mô hình ngôn ngữ lớn, tìm kiếm vector và mô hình sinh ảnh mở ra khả năng xây dựng một trải nghiệm mua sắm mới. Người dùng không chỉ duyệt danh sách sản phẩm, mà có thể trò chuyện với trợ lý AI để nhận gợi ý phối đồ, sử dụng tủ đồ cá nhân để phối cùng sản phẩm trong cửa hàng và thử đồ ảo trên ảnh của chính mình. Đề tài vì vậy hướng tới việc kết hợp thương mại điện tử truyền thống với các thành phần AI để tạo một hệ thống mua sắm thời trang có tính cá nhân hóa và tương tác cao hơn.

\section*{2. Mục tiêu của đề tài}
\phantomsection
\addcontentsline{toc}{section}{2. Mục tiêu của đề tài}

Mục tiêu tổng quát của đề tài là xây dựng một hệ thống thương mại điện tử thời trang tích hợp AI ở mức sản phẩm khả dụng tối thiểu. Các chức năng thương mại điện tử nền tảng được kết hợp với những quy trình AI phục vụ tư vấn và thử đồ. Hệ thống cho phép người dùng xem sản phẩm, thêm sản phẩm vào giỏ hàng, thiết lập hồ sơ AI và quản lý tủ đồ cá nhân. Trên nền tảng đó, trợ lý phối đồ AI (AI Stylist) phân tích yêu cầu ngôn ngữ tự nhiên và đề xuất các phương án phối đồ dựa trên sản phẩm cửa hàng cùng trang phục cá nhân. Chức năng thử đồ ảo (Virtual Try-On -- VTON) sinh ảnh thử đồ từ ảnh người dùng và ảnh trang phục.

\medskip

Về mặt kiến trúc, hệ thống xử lý bất đồng bộ các tác vụ AI có tải tính toán lớn bằng RabbitMQ để tránh làm chậm yêu cầu thương mại điện tử thông thường. LanceDB lưu vector nhúng của sản phẩm và tủ đồ cá nhân để phục vụ truy hồi ngữ nghĩa trong luồng gợi ý. Kết quả VTON được cập nhật tới giao diện người dùng bằng Server-Sent Events (SSE) ngay khi bộ xử lý hoàn tất tác vụ.

\section*{3. Đối tượng sử dụng và phạm vi}
\phantomsection
\addcontentsline{toc}{section}{3. Đối tượng sử dụng và phạm vi}

Đối tượng sử dụng chính của hệ thống là khách hàng mua sắm thời trang trực tuyến. Người dùng có thể xem sản phẩm, tạo hồ sơ cơ thể, tải lên trang phục cá nhân, trò chuyện với AI Stylist để nhận gợi ý phối đồ và thử đồ ảo trước khi thêm vào giỏ hàng.

\medskip

Hệ thống cũng có một công cụ quản trị đơn giản để thêm sản phẩm mới. Khi sản phẩm được tạo trong MedusaJS, hệ thống nghiệp vụ phát sự kiện đồng bộ siêu dữ liệu qua RabbitMQ để dịch vụ AI phân tích và cập nhật vào LanceDB.

\medskip

Phạm vi đề tài dừng ở mức sản phẩm khả dụng tối thiểu về kỹ thuật. Hệ thống đã tích hợp các mô hình AI, hàng đợi thông điệp, cơ sở dữ liệu vector và cơ chế cập nhật kết quả theo thời gian thực, nhưng chưa phải một sản phẩm sẵn sàng vận hành thực tế. Các nội dung như thanh toán trực tuyến, xác thực người dùng đầy đủ, lưu trữ ảnh dùng chung, cơ chế xử lý lại và hàng đợi lỗi vẫn là những hướng phát triển tiếp theo.

\section*{4. Phương pháp thực hiện}
\phantomsection
\addcontentsline{toc}{section}{4. Phương pháp thực hiện}

Quá trình phát triển được chia thành năm giai đoạn. Trước hết, nhóm khảo sát bài toán mua sắm thời trang trực tuyến và xác định hai chức năng AI trọng tâm là tư vấn phối đồ và thử đồ ảo. Từ phạm vi đã xác định, các yêu cầu chức năng, yêu cầu phi chức năng và luồng tương tác chính được phân tích để làm cơ sở thiết kế kiến trúc.

\medskip

Ở giai đoạn thiết kế, hệ thống được phân tách thành giao diện người dùng, hệ thống nghiệp vụ thương mại điện tử và dịch vụ AI. Các luồng có độ trễ thấp được xử lý theo cơ chế yêu cầu--phản hồi; tác vụ sinh ảnh được đưa qua hàng đợi để xử lý bất đồng bộ. Thiết kế dữ liệu cũng tách dữ liệu giao dịch trong PostgreSQL khỏi vector nhúng phục vụ truy hồi ngữ nghĩa trong LanceDB.

\medskip

Ở giai đoạn cài đặt, từng thành phần được phát triển theo hợp đồng dữ liệu đã thiết kế, sau đó tích hợp thành các luồng hoàn chỉnh từ giao diện, hệ thống nghiệp vụ, dịch vụ AI và ngược lại. Việc kiểm thử tập trung vào khả năng trao đổi dữ liệu giữa các thành phần, trạng thái tác vụ bất đồng bộ và tính đúng đắn của kết quả hiển thị. Cuối cùng, nhóm tiến hành thực nghiệm VTON trên cùng một bộ ảnh đầu vào, ghi nhận thời gian xử lý và phân tích sự khác biệt giữa CatVTON cục bộ với FLUX.2 qua dịch vụ đám mây.

\medskip

Mã nguồn của đề tài được quản lý tại:
\begin{center}
\url{https://github.com/doananhhung/E-commerce-AI-stylist}
\end{center}

\clearpage
